% SYNTHETIC TEST UNIT (integration tools) -- deliberate errors are marked "DELIBERATE"
\chapter{Drell--Yan（合成测试）}\label{ch:dy}

\section{因子化}\label{sec:dy-fact}

\begin{theorem}[DY 因子化 \tagA]\label{thm:dy-factorization}
合成测试。
非微扰地位：\cref{asm:intro-bridge}（MA1+MA2）。
% errtype: E-N
\uses{thm:method-template,asm:appN-collins,asm:intro-bridge}
\end{theorem}
\begin{proofsketch}略。\end{proofsketch}

\paragraph{地位与缺口}全阶论证；瓶颈\cref{asm:appN-collins}；具名缺口\cref{asm:appN-brst-nonlocal}（经\cref{thm:method-template}）。

\begin{proposition}[模型命题 \tagP]\label{prop:dy-model}
在标量模型中、$O(g^4)$ 阶，合成测试。（本命题只涉及模型。）
\uses{thm:dy-factorization}
\end{proposition}

\paragraph{地位与缺口}部分：只在模型中证明；瓶颈\cref{asm:appN-collins}；具名缺口\cref{asm:appN-brst-nonlocal}。

% DELIBERATE C-P16/K8 E no status paragraph
\begin{conjecture}[有色末态 \tagC]\label{conj:dy-coloured}
合成测试。（本猜想只涉及部分子态。）
\uses{thm:dy-factorization}
\end{conjecture}

% DELIBERATE C-P16/K8 W status paragraph not immediately after the environment
\begin{proposition}[地位段落滞后 \tagA]\label{prop:dy-late}
合成测试。（本命题只涉及部分子态。）
\uses{thm:method-template}
\end{proposition}

中间插入的一段正文（合成测试）。

\paragraph{地位与缺口}全阶论证；瓶颈\cref{asm:appN-collins}；具名缺口\cref{asm:appN-brst-nonlocal}。

% DELIBERATE C-P15 E the 非微扰地位 line names a technical assumption
\begin{proposition}[非微扰地位行含技术假设 \tagA]\label{prop:dy-npline}
合成测试。
非微扰地位：\cref{asm:intro-bridge}、\cref{asm:appN-collins}。
\uses{thm:dy-factorization,asm:intro-bridge,asm:appN-collins}
\end{proposition}

\paragraph{地位与缺口}全阶论证；瓶颈\cref{asm:appN-collins}；具名缺口\cref{asm:appN-brst-nonlocal}。

% DELIBERATE C-P15 E hadron-level by the master table (column 非微扰地位 = Br) but no 非微扰地位 line
\begin{proposition}[强子层面命题 \tagA]\label{prop:dy-hadron}
合成测试。
\uses{thm:dy-factorization}
\end{proposition}

\paragraph{地位与缺口}全阶论证；瓶颈\cref{asm:appN-collins}；具名缺口\cref{asm:appN-brst-nonlocal}。

% DELIBERATE K10 E forbidden phrasing next to an [A] statement
\begin{proposition}[禁用措辞 \tagA]\label{prop:dy-forbidden}
合成测试。（本命题只涉及部分子态。）
\uses{thm:dy-factorization}
\end{proposition}

\paragraph{地位与缺口}本命题已被证明（措辞违规的合成测试）；瓶颈\cref{asm:appN-collins}；具名缺口\cref{asm:appN-brst-nonlocal}。

% DELIBERATE C-P5 E an [A] statement uses the [C] technical assumption (the named gap as a \uses premise)
\begin{proposition}[使用 C 假设的命题 \tagA]\label{prop:dy-ctech}
合成测试。（本命题只涉及部分子态。）
\uses{thm:dy-factorization,asm:appN-brst-nonlocal}
\end{proposition}

\paragraph{地位与缺口}全阶论证；\cref{asm:appN-brst-nonlocal}。

% DELIBERATE C-P2/K3 E canonical label owned by Ch07 defined in Ch12
\begin{proposition}[放错单元的命题 \tagR]\label{prop:dy-misplaced}
（$R_{\mathrm{PT}}$：在重整化微扰论中对红外正规化的部分子态严格。）合成测试。（本命题只涉及部分子态。）
\uses{def:dis-kinematics}
\end{proposition}

% DELIBERATE C-P13 W retired pending-verification comment, withdrawn slot label; K2 E foreign equation label
% PENDING-VERIFY: synthetic marker
另见\cref{thm:appN-oneloop}与\cref{eq:dis-foreign}。

% DELIBERATE figsync NOSOURCE (block not in the figure sources)
%%% BEGIN fig:synth-nosource
\begin{figure}[htbp]
\centering
\fbox{合成测试图丁}
\caption{合成测试图丁（图源中没有）。}\label{fig:synth-nosource}
\end{figure}
%%% END fig:synth-nosource

% DELIBERATE figsync DUPLICATE / C-P12 / K1 the block fig:synth-same is also in ch07
%%% BEGIN fig:synth-same
\begin{figure}[htbp]
\centering
\begin{tikzpicture}[fd]
\draw[quark] (0,0) -- (2,0);
\end{tikzpicture}
\caption{合成测试图甲（与单元中的副本一致）。}\label{fig:synth-same}
\end{figure}
%%% END fig:synth-same
